\section{Implementation and Deployment}

\subsection{Technology Stack}
The implementation is modular. Python components build the corpus, index, graph, and evaluation outputs. FastAPI exposes the runtime assistant, Qdrant provides hybrid retrieval, Neo4j stores legal structure, and the Next.js frontend uses authenticated server-side proxy routes.

\begin{table}[H]
  \centering
  \caption{Technology Stack and System Modules}
  \label{tab:tech-stack-draft}
  \begin{tabularx}{\linewidth}{@{}L{0.22\linewidth}L{0.25\linewidth}Y@{}}
    \toprule
    Layer & Technology & System role \\
    \midrule
    Corpus pipeline & Python & Validate, chunk, enrich, and export legal evidence. \\
    Embeddings & Sentence Transformers & Build dense Vietnamese SBERT vectors. \\
    Vector search & Qdrant & Dense and sparse hybrid retrieval with payload metadata. \\
    Graph store & Neo4j & Legal hierarchy, references, taxonomy, and expansion. \\
    Backend API & FastAPI & Chat, health, auth, admin, and conversation routes. \\
    Frontend & Next.js & Chat interface, admin pages, and authenticated backend proxy. \\
    Evaluation & Python & Retrieval ablation and end-to-end metrics. \\
    \bottomrule
  \end{tabularx}
\end{table}

\subsection{Offline Index and Graph Build}
The offline index build resolves enriched legal chunks and creates dense text, sparse BM25-style lexical vectors, Qdrant payload records, keyword payload indexes, and an index manifest. The manifest records 1,556 indexed chunks, six document groups, \texttt{keepitreal/vietnamese-sbert}, 768-dimensional dense vectors, a sparse vector channel, and the active Qdrant collection. Keeping this manifest makes evaluation reproducible without rebuilding the corpus for every experiment.

The graph build creates evidence-chunk nodes, legal document/article/clause/point/appendix nodes, taxonomy nodes, concept nodes, source links, reference edges, guidance/detail edges, and normative hierarchy edges. The graph is more expensive to prepare than a pure vector index, but it gives the retriever explicit paths for cross-references, implementation guidance, and hierarchy-aware expansion.

\subsection{Runtime Backend and \texttt{/chat} Flow}
The FastAPI application configures CORS from runtime settings and includes health, auth, conversation, chat, and admin routers. Shared dependencies lazily construct the \texttt{HybridRetriever} from the configured index path and reranker settings, and they enforce bearer-token authentication before protected routes are executed.

The \texttt{/chat} route acts as the integration boundary between the application layer and the RAG pipeline. It authenticates the user, resolves the conversation, retrieves evidence with configurable \texttt{topK} and \texttt{prefetchLimit}, and reduces retrieved contexts under count, character, and token budgets. When \texttt{retrieveOnly} is enabled, the route returns the selected evidence directly; otherwise, it passes the user message and bounded contexts to the grounded answer generator.

The route also stores the assistant response with extracted legal bases and evidence quotes, while \texttt{ChatTraceService} records latency, routed intent, retrieved hits, selected contexts, citations, insufficient-context status, and errors. This trace layer is useful for evaluating failures because it connects a final answer to the retrieval and context-selection decisions that produced it.

If retrieval returns no usable context, the route stores an insufficient-context answer and records a trace instead of generating unsupported legal advice. This behavior mirrors the thesis safety requirement: answer construction is downstream of retrieval evidence, not independent from it.

\subsection{Answering and Safety Implementation}
Grounded answer construction is implemented in the answer generation module listed in Appendix~A. It selects answer contexts, checks the query against the scope guard, then either uses the configured LLM provider or deterministic extractive synthesis. If an LLM call fails or produces an invalid grounded answer, the code can fall back to deterministic extractive output.

Citation validation is implemented in the validation module. It verifies retrieved legal bases, unsupported article numbers, higher-rank source ordering, and insufficient-context wording. Scope refusal rules handle unsupported topic families before answer generation proceeds.

\subsection{Configuration}
Runtime configuration is centralized in a settings object loaded from environment variables and the repository \texttt{.env} file. It covers authentication, CORS, Qdrant connection details, index path, reranker settings, article sibling contexts, embedding provider, query router, Neo4j connection, legal graph expansion settings, model provider settings, and benchmark settings.

This allows the same codebase to run locally, in evaluation mode, or in hosted deployment. The deterministic extractive provider used in the final evaluation avoids provider variability and makes benchmark reproduction easier.

\subsection{Frontend Proxy and Deployment Architecture}
The frontend uses a Next.js server-side proxy instead of calling the backend directly from the browser. The server API helper defines \texttt{BACKEND\_URL} resolution, the \texttt{auth\_token} cookie name, secure cookie options, and bearer authorization headers. The proxy module reads the cookie from the incoming \texttt{NextRequest}, rejects unauthenticated requests, forwards JSON or streaming requests to the FastAPI backend, and preserves the backend status or stream headers.

In deployment, the Next.js frontend can run on Vercel with \texttt{BACKEND\_URL} pointing to the FastAPI service. The backend can run on Render with environment variables for authentication, Qdrant, Neo4j, and optional provider credentials. Qdrant may run as Qdrant Cloud, while Neo4j can be provided by a managed instance or containerized service.

\thesisfigure{deployment_architecture.pdf}{Deployment Architecture with Vercel frontend and Render backend.}{fig:deployment}

\subsection{Evaluation Reproducibility}
The evaluation path is deterministic. Retrieval ablation runs hybrid-only and graph-augmented modes over the frozen benchmark and aggregates recall, required citation coverage, forbidden citation violation rate, and retrieval pass rate. The end-to-end run enables graph retrieval, generates grounded answers, validates answer quality and citations, and aggregates category, topic, and difficulty metrics. Appendix~A lists the fixed inputs and structured outputs used for reproduction.

The final deterministic run is extractive, with top-k set to 10, prefetch limit set to 24, max answer contexts set to 8, and graph expansion used in 89 queries.
